% ============================================================================
% Section 9.6 — Finding Probabilities of Compound Events
% CCSS 7.SP.C.8b
% ============================================================================
\section{Finding Probabilities of Compound Events}

\begin{stepGoal}
\begin{itemize}
    \item Calculate the probability of a compound event from a sample space
    \item Apply compound probability to everyday situations
\end{itemize}
\end{stepGoal}

\begin{stepVocabBox}
    \vocabItem{Compound Probability}{The probability that \textbf{two or more events} happen together.}
    \vocabItem{Favorable Outcomes}{Outcomes in the sample space that match your event.}
\end{stepVocabBox}

\begin{stepsCard}[How to Find the Probability of a Compound Event]
    \stepItem{Build the \textbf{sample space} (list, table, or tree diagram).}
    \stepItem{Count the \textbf{total outcomes} in the sample space.}
    \stepItem{Count the \textbf{favorable outcomes} that match the event.}
    \stepItem{Divide: $P = \dfrac{\text{favorable outcomes}}{\text{total outcomes}}$ and simplify.}
\end{stepsCard}

% --- Worked Example 1 ---
\begin{stepExample}{You roll two dice. What is the probability that the sum is $7$?}
    \stepShow{1} Sample space (table of sums):

    \begin{center}
    \small
    \begin{tabular}{c|cccccc}
    $+$ & $1$ & $2$ & $3$ & $4$ & $5$ & $6$ \\ \hline
    $1$ & $2$ & $3$ & $4$ & $5$ & $6$ & $7$ \\
    $2$ & $3$ & $4$ & $5$ & $6$ & $7$ & $8$ \\
    $3$ & $4$ & $5$ & $6$ & $7$ & $8$ & $9$ \\
    $4$ & $5$ & $6$ & $7$ & $8$ & $9$ & $10$ \\
    $5$ & $6$ & $7$ & $8$ & $9$ & $10$ & $11$ \\
    $6$ & $7$ & $8$ & $9$ & $10$ & $11$ & $12$ \\
    \end{tabular}
    \end{center}

    \stepShow{2} Total outcomes: $6 \times 6 = 36$.

    \stepShow{3} Sums of $7$: $(1,6), (2,5), (3,4), (4,3), (5,2), (6,1)$ --- that is $6$ favorable.

    \stepShow{4} $P(\text{sum} = 7) = \dfrac{6}{36} = \dfrac{1}{6}$
    \stepResult{$P(\text{sum} = 7) = \dfrac{1}{6}$}
\end{stepExample}

% --- Worked Example 2 ---
\begin{stepExample}{You flip two coins. What is the probability of getting at least one head?}
    \stepShow{1} Sample space: $HH, HT, TH, TT$.

    \stepShow{2} Total outcomes: $4$.

    \stepShow{3} Favorable (at least one $H$): $HH, HT, TH$ --- that is $3$.

    \stepShow{4} $P(\text{at least one } H) = \dfrac{3}{4}$
    \stepResult{$P(\text{at least one head}) = \dfrac{3}{4}$}
\end{stepExample}

\mascotStepTip{``At least one'' means one or more. It's often easier to count the outcomes that do NOT match and subtract from the total.}

% ============================================================================
% PRACTICE
% ============================================================================
\newpage
\begin{stepPractice}{Compound Probability Practice}
    \resetProblems

    \stepPracticeHeader[step1Color]{Build and Count the Sample Space}
    \prob You spin a spinner ($1, 2, 3$) and flip a coin. How many total outcomes? \answerBlank[2cm]
    \answer{$3 \times 2 = 6$}

    \stepPracticeHeader[step2Color]{Find Favorable Outcomes}
    \prob Using the spinner and coin above, how many outcomes have an even number \textbf{and} heads? \answerBlank[2cm]
    \answer{$1$ (only $2H$)}

    \stepPracticeHeader[step3Color]{Calculate the Probability}
    \prob What is $P(\text{even and heads})$ from the problem above? \answerBlank[2cm]
    \answer{$\frac{1}{6}$}

    \prob You roll two dice. What is the probability of getting a sum of $2$? \answerBlank[2cm]
    \answer{$\frac{1}{36}$ (only $(1,1)$)}

    \prob You roll two dice. What is the probability that both show the same number (doubles)? \answerBlank[2cm]
    \answer{$\frac{6}{36} = \frac{1}{6}$}

    \wordProblem{A bag has $2$ red and $3$ blue marbles. You pick one marble, put it back, and pick again. What is the probability of picking red both times?}{~}
    \answer{$P = \frac{2}{5} \times \frac{2}{5} = \frac{4}{25}$}
\end{stepPractice}
